%% ma: L11-B11-0006
%% lop: 11
%% chuong: 4
%% bai: 11
%% dang: 11.11.3
%% loai: TLN
%% muc: VD
%% dap_an: "1,33"
%% nguon: "de-cuong-giua-ky-1-toan-11-nam-2026-2027-truong-thpt-hoang-van-thu-ha-noi-2-8.pdf (M08, Đề cương giữa kì 1 Toán 11 THPT Hoàng Văn Thụ, Hà Nội), Đề minh họa, Phần III Câu 4"
%% kien_thuc: "Bài 10 (giao tuyến), Bài 11 (định lí về giao tuyến song song), định lí Thalès (lớp 8)"
%% trang_thai: chinh-thuc
%% ghi_chu: "Đã duyệt 29/10/2026."
\begin{ex}
Cho tứ diện $ABCD$, gọi $I$ và $J$ lần lượt là trung điểm của $AD$ và $AC$, $G$ là trọng tâm của tam giác $BCD$, mặt phẳng $(GIJ)$ cắt $BC$, $BD$ lần lượt tại $M$ và $N$. Đặt $MN=x\cdot IJ$. Tìm $x$ (kết quả làm tròn đến hàng thập phân thứ hai).
\shortans{1,33}
\loigiai{$IJ\parallel CD$ và $IJ=\dfrac12CD$. Hai mặt phẳng $(GIJ)$, $(BCD)$ chứa $IJ\parallel CD$ và có điểm chung $G$, nên giao tuyến $MN$ qua $G$ và song song với $CD$.

Gọi $E$ là trung điểm $CD$; $\dfrac{BG}{BE}=\dfrac23$ nên $MN=\dfrac23CD$. Vậy $x=\dfrac{2/3}{1/2}=\dfrac43\approx1{,}33$.}
\end{ex}
